%% Acknowledgements
\begin{preface}

The starting point for this thesis was my multidisciplinary interest, which made me want to apply statistical rigour to a real-world problem. This is the result of that attempt.

I began the thesis work by using simulations to explore and visualise intricacies of the problem of estimating the causal effect of social media use on well-being. After the first two months, I realised the need for more statistical rigour, and made a major change of direction. This led me to the topic for the main body of the thesis: developing statistical inference for network models.

The result is a bidirectional story. Complex simulations apply a bottom-up approach, showing the difficulties of the problem. Top-down, statistically rigorous models point towards a resolution, but do not yet have the sophistication required for a full answer. I have enjoyed working with the problem in this way, and hope that my contribution provides a step towards closing the gap in the future.

I want to thank Tobias Dahl for his energy, optimism, and encouragement throughout the process. I would also like to thank my internal supervisor Ingrid Hobæk Haff for help with the formal requirements of the programme. Special thanks go to Steffen Grønneberg for tremendous support during the most uncertain time of the work, and for helping me find the way forward.

Finally, I want to thank my family and friends for all the interest they have shown in this project, and for their feedback on drafts.

\vspace{3cm}
\noindent Oslo, May 2026

\vspace{1cm}
\noindent Aslak Hellevik

\end{preface}
